% !TeX root = ../papers/nc-p3.tex
\section{半単純対称性の四位置での一致と三つの完全分類}
\label{sec:sl2-window-profiles}
$L_d=\operatorname{Sym}^{d-1}k^2$ を $\mathrm{SL}_2$ の $d$ 次元既約表現とする。
\begin{theorem}\label{thm:sl2-window-profiles}
標準AS窓の射影安定群への非自明な $\mathrm{SL}_2$ 準同型を取る。四位置の矢表現は全て同じ型であり、
各関係表現も全位置で同じになる。可能な型と、最後の三行で印を忘れた全正則窓は次の通りである。
\[
\begin{array}{c|c|l}
 V_i & R_i & \text{窓}\\\hline
 L_2+2L_1 & 2L_2+2L_1 & \text{下記の二次元係数表示}\\
 L_3+L_1 & 2L_3 & \text{多項式環または }U(\mathfrak{sl}_2)\text{ のPBW斉次化}\\
 2L_2 & L_3+3L_1 & \text{多項式環または下記の非自明な単冪型}\\
 L_4 & L_5+L_1 & \text{多項式環}
\end{array}
\]
三つの同型類の最小周期は全て1で、反対階数は多項式環6、斉次包絡環8、単冪型12である。
\end{theorem}
\begin{proof}[全表現の制限]
単連結性により四射影表現を線形に持ち上げ、最上文字は自明である。
各四次元表現は4の分割、各六次元表現は6の分割であり、相補二位置の関係表現は双対、従って同型である。
$L_i\otimes L_j=L_{i+j-1}+L_{i+j-3}+\cdots+L_{|i-j|+1}$ から各隣接積と対辺の積の
両方に入る六次元部分表現を列挙する。極大トーラスのEuler四文字はWeyl反転で不変だから
$\{m,-m,n,-n\}$ である。$p_1$ の公式で各文字の絶対値は12以下となる。
\path{computations/verify_sl2_window_profiles.py} は整数値を一つの素数で還元した不一致だけで
この表と自明作用以外を全て落とす。$L_4$ の場合 $L_5+L_1=\bigwedge^2 L_4$ は積で一意な部分表現なので、
全辺の関係は四位置を同一視した通常の交換子である。
\end{proof}

\begin{proof}[$L_3+L_1$]
矢をスカラー $t_i$ と標準三次元直交空間の座標 $x$ に分ける。各辺に現れる三次元型の
三つのコピーは $t_i x,x t_{i+1},x\times x$ で、最後は三次元交代形式による包含を表す。
隣接する二スカラーの成分はない。四つとも $x$ の部分は第一切断で交代積二つの内積、
第二切断でも巡回した同じ形の倍数である。独立な三つの計量の対積で展開するとこの二式は比例せず、
四つとも $x$ の部分は零である。残る一スカラー成分を
$a_i t_i\epsilon(x_{i+1},x_{i+2},x_{i+3})$、二つの反対スカラーの成分を順に
$u t_0t_2(x_1,x_3),v t_1t_3(x_2,x_0)$ と書く（添字順のテンソルを表す）。
上述の三コピーで第一切断の重複度行列は
\[
 \begin{pmatrix}u&0&a_0\\0&v&a_1\\a_2&a_3&0\end{pmatrix}.
\]
他の辺では巡回する。必要な階数2から $v a_0a_2+u a_1a_3=0$。
$u=0$ なら一文字全階数から $a_0a_2\ne0$、従って $v=0$ となる。
このとき全 $a_i\ne0$ で、列空間は $x\times x$ と
$a_i t_i x+a_{i+1}x t_{i+1}$ を与える。位置別の $t_i$ の倍率で通常の交換子表示になる。
$uv\ne0$ なら $a_{i+2}=0$ のとき、前辺に $x t_i$、次辺に $t_i x$ が純関係として続くので
正則でない。全 $a_i$ を1にして共通倍率を調整すると $u=1,v=-1$。
$(-1)^i t_i$ を新しい中心文字にすれば残る三式は三次元交代積によるLie括弧の
斉次化である。包絡環のPBW定理または定理\ref{thm:regular-normal-linear}で全次数正則となる。
\end{proof}

\begin{proof}[$2L_2$]
$V_i=U_i\otimes L_2$, $\dim U_i=2$ とする。
四つの標準表現の不変形式は $[01][23],[12][30]$ が張る。ここで $[ij]$ は交代二形式である。
第一切断の $L_3$ の重複度は1なので、第二の不変形式の係数四テンソルは
$01|23$ で階数1である。第二切断から第一の係数も $12|30$ で階数1である。
従って非零二行列 $r_i$ を使い、全窓のテンソルが
\[
   A\, r_1 r_3[01][23]+B\,r_0 r_2[12][30],\qquad AB\ne0
\]
となる。$r_i$ は位置 $i,i+1$ の重複度添字に置く。
$J=\left(\begin{smallmatrix}0&1\\-1&0\end{smallmatrix}\right)$ として
$r_i=JF_i^{\mathsf t}$ と書くと、$F_i$ の基底変更は倍率を除き有向サイクルの写像の基底変更である。
自明表現成分の階数は各切断で3でなければならない。第一で
$2A r_1r_3+B r_0r_2$、第二で $A r_1r_3+2B r_0r_2$ を当該順で切る。
各々一方の項は階数1、他方は対辺の二行列の階数の積なので、二つの対辺がともに特異にはならない。

特異な辺が高々一つなら、巡回と基底変更後に $(F_0,F_1,F_2,F_3)=(I,I,I,M)$ とする。
二つの四次行列式は
\[
 -8A^3\det M(2A\det M+B\operatorname{tr}M),\qquad
 -8B^3(A\operatorname{tr}M+2B).
\]
$M$ が可逆なら $(\operatorname{tr}M)^2=4\det M$。
倍率を吸収して $A=1,B=-1$、$M=I$ または非自明な単冪Jordan行列とできる。
前者のテンソルは四つの積文字の交代体積である。両者とも標準表現の対角トーラスで
二平面ずつに分ければ純次数の二形式は $r_i$ の非零倍で、全て可逆である。
従って定理\ref{thm:split-inverse-regularity}が使える。
$M$ の行列式1の四乗根を全辺に置く表示に移せるので、どちらも一巡回で符号だけが変わり周期1である。
$M$ が特異なら跡は非零なので $M=\operatorname{diag}(1,0)$, $B=-A/2$ とでき、
始点1の四次元数39で除外される。

特異な二辺は隣接するので $F_0=F_1=I,F_2=\operatorname{diag}(1,0)$ とする。
二つの頂点空間で核と相手の像が一致するかどうかに従って
$F_3=E_{00},E_{11},E_{01},E_{10}$ の四場合になる。中央二つは隣接階数が5以下である。
他の二場合では $A=1,B$ 独立の関数体上で始点1の四次元数52となる。
これは実際の関係像の階数が6の全特殊化でも下界である。上記コードは四文字の全語に
二次関係を三通りに挿入し、標準表現のトーラス次数ごとに正確に掃き出す。
反対階数6,8,12も表示自体から直接得られるので、印を忘れても三種は区別される。
\end{proof}

\begin{proposition}\label{prop:one-standard-sl2-slice}
$L_2+2L_1$ の場合、二次元の自明文字を $U_i$、標準文字を $S_i$ とする。
テンソルは $P\in\bigotimes U_i$ と、二つの $S$ に交代形式、その補二位置に二行列を置く六成分で尽くされる。
$U_i,U_{i+1}$ を使用する境界二行列を $b_i$ とする。正則性には全 $b_i\ne0$ が必要である。
全 $b_i$ が可逆なら所定の二隣接表現階数を満たす全係数が標準Koszul左右AS正則である。
\end{proposition}
\begin{proof}
辺に $L_3$ がないので、四標準文字の成分は四つの各隣接対で交代的であり零である。
奇数個の標準文字の成分もない。$U$ を全て消す商では各辺の関係の像は標準二文字の
交代関係の倍数だけである。$b_i=0$ なら対辺のこの関係も零になる。
実際この射影には二つの標準文字を含む当該成分と、既に消えた四標準文字成分だけが寄与する。
四文字ごとに少なくとも一つの関係のない境界があり、間のブロックは少なくとも二つの
対称語を残すので指数増大となる。

全 $b_i$ 可逆の場合、各 $S_i$ を同率で拡大して射影極限を取り、$P$ だけを零にする。
混合対切断の標準表現成分は変わらず重複度2である。自明部分では、
純 $UU$ 対と $SS$ の交代成分との間の非零列・行だけが残り、その和の階数は2となる。
極限は $2+2$ の二平面分解を持ち、純次数関係は $b_i$ と $S_i$ の交代形式である。
全て可逆なので前記定理で正則であり、命題\ref{prop:unrestricted-good-reduction}で元の窓も正則となる。
\end{proof}
